\section{Cái chương này đóng được, và cái nó mở ra}
\label{sec:ch10-con-lai-gi}

\index{thiên kiến quy nạp!chương này đóng được gì}%
Chương~\ref{ch:chi-phi-va-song-song} để lại ba khoảng trống của kiến trúc
Transformer, và chương~\ref{ch:tien-huan-luyen} đã lấy khoảng trống thứ ba, đói
dữ liệu. Chương này lấy khoảng trống thứ hai, và
chương~\ref{ch:chi-phi-va-song-song} gọi tên nó đầy đủ hơn tôi vừa gọi: không
có \tn{thiên kiến quy nạp}{inductive bias} \emph{về vị trí}. Chữ cuối đáng giữ,
vì hai ràng buộc của mục~\ref{sec:ch10-hai-rang-buoc} đều là phát biểu về không
gian.

Cái nó đóng được thì gọn, và gọn hơn tôi tưởng lúc bắt đầu.

Thiên kiến quy nạp của một mạng \tn{tích chập}{convolution} là hai ràng buộc trên một
\tn{ma trận trọng số}{weight matrix}, và cả hai đếm được.
\tn{Tính cục bộ}{locality} chia số tham số cho 109.75 ở tầng đầu của mạng
mục~\ref{sec:ch10-anh-that}; \tn{chia sẻ trọng số}{weight sharing} chia tiếp
cho đúng 1024, tức đúng số vị trí không gian. Hai ràng buộc không ngang nhau, và cái lớn hơn là cái ít
được nói tới hơn.

Cái chúng mua không phải kích thước. Bộ trích đặc trưng của LeNet-5 nặng 1,716
tham số trong tổng 60,000, tức chia sẻ trọng số làm nhỏ đi đúng phần nó chạm
tới và không làm nhỏ mạng. Cái chúng mua là dữ liệu: 28.4 lần, và một MLP với
23.71 lần tham số vẫn cần 11.0 lần. Đo trên CIFAR-10 ở quy mô một laptop.

\subsection*{Hai con số quay về, và một tính chất không}

\index{thiên kiến quy nạp!hai phép dựng lại đều khớp}%
Chương~\ref{ch:tien-huan-luyen} dựng lại bốn bài và ba bài không khớp. Chương
này dựng lại hai bài và cả hai khớp tới đơn vị: 64,660 kết nối với 9,760 tham
số ở bài 1989, và 340,908 với 60,000 ở LeNet-5, trong đó 60,000 là tổng đúng
chứ không phải một phép làm tròn.

Tôi in chỗ này ra vì nó là chiều ngược của bài học chương trước. Ở đó kết luận
là một con số được trích lại đủ nhiều thì không còn ai dựng lại. Ở đây phép
dựng lại chạy và không tìm ra gì, và điều đó cũng đáng biết: một quy tắc chỉ
chạy khi người ta đoán nó sẽ hỏng thì không phải quy tắc, nó là một linh cảm.

Cái \emph{không} quay về là một tính chất chứ không phải một con số. Mạng tích
chập được cả ngành mô tả là \tn{bất biến}{invariance} với phép dịch. Đo ra thì
tầng tích chập \tn{đẳng biến}{equivariance} tuyệt đối - đúng số không, không
phải số nhỏ - và phép \tn{giảm mẫu}{subsampling} làm
mất tính chất ấy ở đúng một nửa số phép dịch, chứ không làm yếu nó đi ở mọi
phép dịch. Mạng ba tầng pooling trên ảnh 32 pixel đẳng biến ở 4 trong 32 phép
dịch ngang.

Và nó không phải chuyện của max pooling. Ba cách giảm mẫu, kể cả một tầng conv
stride 2 không có phép max nào, cho cùng một dạng chẵn lẻ. Bài 1989 đã viết
\enquote{some position information is eliminated} về chính phép giảm mẫu của
nó; cái mất ba mươi năm là đo xem bao nhiêu.

\subsection*{Một phép đo tôi làm sai, và cách nó lộ ra}

\index{thiên kiến quy nạp!phép đo đẳng biến sai ở bản đầu}%
Bản đầu của phép đo đẳng biến so \(f(T_\delta x)\) với \(T_\delta f(x)\) ở đúng
cùng một \(\delta\), kể cả khi \(f\) giảm mẫu. Với một tầng giảm mẫu hai lần
thì so như vậy là sai đại lượng: dịch đầu vào 2 phải cho đầu ra dịch 1.

Bảng chạy ra sai lệch lớn ở \emph{mọi} phép dịch, kể cả các phép dịch chẵn mà
tầng ấy đúng tuyệt đối. Đọc bảng ấy thì kết luận là \enquote{giảm mẫu phá đẳng
biến ở mọi phép dịch}, mạnh hơn hẳn và sai.

Không có gì trong build hay trong gate bắt được nó. Cả hai bảng đều là số thật,
chạy thật. Cái bắt được nó là đem số đo so với chuyện mình đoán trước, thấy
lệch, rồi đi tìm lý do thay vì in ra - và ở đây tôi đoán trước rằng phép dịch
chẵn phải cho số không, nên khi nó không cho thì có một chỗ để kéo.

Cùng loại với chuyện ở mục~\ref{sec:ch10-anh-that}, nơi tôi suýt viết rằng
khoảng cách hẹp lại theo dữ liệu trong khi bảng nói nó rộng ra. Cả hai lần, cái
cứu là có một dự đoán trước khi nhìn số.

\subsection*{Cái nó mở ra}

\index{thiên kiến quy nạp!chỗ hai chương sau bắt đầu}%
Hai chỗ, và chúng đi hai hướng.

Thứ nhất, chương~\ref{ch:vit}, và chương này vừa dựng xong đúng cái nền nó cần.
Mục~\ref{sec:ch10-attention-lam-duoc} cho thấy khác biệt giữa hai kiến trúc
không nằm ở sức biểu diễn: self-attention với \(K^2\) head biểu diễn được tích
chập \(K \times K\), và điều đó chứng minh được chứ không phải phỏng đoán. Khác
biệt nằm ở chỗ cái nào là mặc định. Tích chập \emph{bắt buộc} cục bộ và dùng
chung trọng số; self-attention thì có thể trở thành như vậy nhưng phải học, từ
dữ liệu.

ViT là chuyện gì xảy ra khi bạn chọn vế thứ hai và có đủ dữ liệu để trả. Bài
nói ra ở mục 3.2 rằng chỉ có đúng hai chỗ trong cả mô hình mà thông tin về cấu
trúc hai chiều được đưa vào bằng tay, và một trong hai chỉ là phép cắt ảnh
thành patch~\autocite{dosovitskiy2021vit}. Bảng của chương ấy là bảng của mục
này với một Transformer thêm vào, và bội số 28.4 lần ở đây là lý do bảng ấy có
hai nửa.

Thứ hai, chương~\ref{ch:chi-phi-bac-hai}. Cả bốn hướng của
mục~\ref{sec:ch10-truoc-vit} đều vấp vào chi phí bậc hai trước khi vấp vào bất
cứ chuyện gì khác, và câu của Image Transformer nói chỗ ấy gọn nhất: một ảnh
32x32 là 3072 vị trí, và attention toàn cục trên nó đã không chạy nổi từ 2018.
ViT không giải quyết chuyện đó, nó đi vòng, bằng cách coi một patch 16x16 là
một token. Chương~\ref{ch:chi-phi-bac-hai} là chương quay lại giải quyết thật.

Còn một chỗ chương này không đóng và cũng không giao cho ai. Bảng ở
mục~\ref{sec:ch10-anh-that} đo CNN cạnh MLP, tức đo hai ràng buộc cùng lúc. Tôi
không tách được phần của tính cục bộ khỏi phần của chia sẻ trọng số bằng bảng
ấy, vì mô hình ở giữa - cục bộ mà không dùng chung - là thứ PyTorch không có
sẵn một tầng cho. Bài tập tier hai cuối chương yêu cầu dựng nó, và tôi ghi ở
đây rằng con số ấy chưa ai đo chứ không ghi rằng nó là bài tập.
